\clearpage
\section{Renaming the complete unary fragment}
\label{sec:copying-unary-renaming}
\termdef{copying-unary-fragment}
The serial compatibility results now yield a full constructor theorem for the
original unary translation, including its final single-use distributors. Every
finite chain of negation, reference creation, and dereferencing over an atomic
term respects bijective variable and atom renaming. In particular this covers
the original $\mathtt{ref}\ 1$ and $!x$ subterms of Proposition~12.2, as well as
dereferencing a named atom. It is a theorem about the original translated graphs;
it does not assert an allocation or read transition of the unfinished machine.

\subsection{Decomposing a typed environment permutation}
An environment consists of $k$ term positions followed by $h$ atom positions.
Embed a variable index $i$ at position $i$, and an atom index $a$ at position
$k+a$. Looking up their wire types gives term and atom, respectively. Conversely,
a position has term type exactly when it is less than $k$, and atom type exactly
when it is at least $k$. Split the appended boundary list at $k$ to prove these
two characterisations. They work with either block empty.

\lean{Copying/EnvironmentPermutation}{8}{8}{EnvironmentPermutation.termIndex}
\lean{Copying/EnvironmentPermutation}{9}{9}{EnvironmentPermutation.atomIndex}
\lean{Copying/EnvironmentPermutation}{11}{12}{EnvironmentPermutation.term_label}
\lean{Copying/EnvironmentPermutation}{14}{15}{EnvironmentPermutation.atom_label}
\lean{Copying/EnvironmentPermutation}{18}{22}{EnvironmentPermutation.term_iff}
\lean{Copying/EnvironmentPermutation}{25}{33}{EnvironmentPermutation.atom_iff}

Let $\pi$ be any bijection of environment positions preserving wire types.
Its variable component sends $i$ to the numeric value of $\pi(i)$; type
preservation proves that value is less than $k$. Its atom component sends $a$
to the numeric value of $\pi(k+a)$ minus $k$. The result is less than $h$ by the
environment bound. Type preservation proves $\pi(k+a)\geq k$, so embedding the
new atom index reconstructs $\pi(k+a)$ exactly. The variable reconstruction is
definitional.

\lean{Copying/EnvironmentPermutation}{35}{36}{EnvironmentPermutation.termMap}
\lean{Copying/EnvironmentPermutation}{38}{40}{EnvironmentPermutation.atomMap}
\lean{Copying/EnvironmentPermutation}{42}{42}{EnvironmentPermutation.term_action}
\lean{Copying/EnvironmentPermutation}{45}{49}{EnvironmentPermutation.atom_action}

The inverse of $\pi$ preserves types as well: apply type preservation at
an inverse image and then the right-inverse equation. Extract the two maps from
this inverse. To prove the extracted variable maps inverse, embed their indices
in the environment and use the original inverse equation; equality of numeric
values suffices for finite indices. For atoms, reconstruct both embedded images
and cancel the common offset $k$. These arguments package independent resource
bijections $\rho$ and $\alpha$ from an arbitrary typed environment bijection.
No premise says in advance that the two sorts are preserved separately.

\lean{Copying/EnvironmentPermutation}{52}{55}{EnvironmentPermutation.inverse_typed}
\lean{Copying/EnvironmentPermutation}{57}{70}{EnvironmentPermutation.terms}
\lean{Copying/EnvironmentPermutation}{72}{91}{EnvironmentPermutation.atoms}

\subsection{Exact actions through the original distributors}
The tensor boundary action on a left input is the left component's input map
followed by the left block injection. The right input, left output and right
output laws are analogous. Each follows from the existing parallel-insertion
commutation theorem, specialised to the relevant component and boundary port,
then the extracted input or output equation. These laws explicitly retain the
variable-before-atom ordering, including the offset of the right block.

\lean{Copying/TensorBoundaryActions}{9}{17}{TensorBoundaryActions.tensor_input_left}
\lean{Copying/TensorBoundaryActions}{19}{27}{TensorBoundaryActions.tensor_input_right}
\lean{Copying/TensorBoundaryActions}{29}{37}{TensorBoundaryActions.tensor_output_left}
\lean{Copying/TensorBoundaryActions}{39}{47}{TensorBoundaryActions.tensor_output_right}

\termdef{copying-single-use-typed-renaming}
For separate resource bijections $\rho$ and $\alpha$, form the tensor of the
original one-use distributor automorphisms, and transport its input interface
through $k\cdot1=k$ and $h\cdot1=h$. This gives an automorphism of the actual
\texttt{singleUse} translation graph. Its term input and output actions are
$i\mapsto\rho(i)$, while its atom actions are $k+a\mapsto k+\alpha(a)$.
For the input laws, use the original use-major position $i+k\cdot0=i$,
then the tensor input law and equality transport. The output laws use the
original distributor output formula and tensor output law. These proofs do not
simplify the single-use graph to an identity wire.

\lean{Copying/SingleUseBoundaryRenaming}{8}{11}{SingleUseBoundaryRenaming.iso}
\lean{Copying/SingleUseBoundaryRenaming}{13}{23}{SingleUseBoundaryRenaming.input_term}
\lean{Copying/SingleUseBoundaryRenaming}{25}{38}{SingleUseBoundaryRenaming.input_atom}
\lean{Copying/SingleUseBoundaryRenaming}{40}{48}{SingleUseBoundaryRenaming.output_term}
\lean{Copying/SingleUseBoundaryRenaming}{50}{58}{SingleUseBoundaryRenaming.output_atom}

Apply this construction to the extracted pair $(\rho,\alpha)$ of any typed
permutation $\pi$. Every environment position is in the variable or atom block;
in each case use its reconstruction equation and the corresponding block action.
The resulting original single-use graph therefore realises precisely $\pi$
on both its inputs and its outputs. This is the critical extension theorem:
no separate distributor-matching premise is required from its caller.

\lean{Copying/SingleUseBoundaryRenaming}{68}{69}{SingleUseBoundaryRenaming.ofTyped}
\lean{Copying/SingleUseBoundaryRenaming}{72}{82}{SingleUseBoundaryRenaming.ofTyped_input}
\lean{Copying/SingleUseBoundaryRenaming}{85}{95}{SingleUseBoundaryRenaming.ofTyped_output}

\subsection{The complete unary constructor theorem}
\termdef{copying-unary-boundary-congruence}
Let $f:G\to H$ be any typed boundary isomorphism between operands with one
term input and the same environment interface. The original unary clause has
three serial components: the unary operation edge, the operand, and the
single-use environment distributor. Its operation edge has one output and the
operand has one input. Both associated index maps are necessarily the identity
on a singleton, so the operation's identity isomorphism glues with $f$.
This constructs the prefix isomorphism with no extra hypothesis.

\lean{Copying/UnaryBoundaryRenaming}{11}{24}{UnaryBoundaryRenaming.unary_prefix_compatible}
\lean{Copying/UnaryBoundaryRenaming}{26}{29}{UnaryBoundaryRenaming.unaryPrefixIso}

The output map $f_O$ preserves boundary wire types by the extracted output
label law. Apply the preceding single-use extension theorem to this very
bijection, obtaining the forest automorphism. The prefix's output action is
$f_O$, by the serial output law; the forest's input action is $f_O$, by the
extension theorem. Thus their compatibility is proved. Applying serial gluing
again gives an isomorphism between the two literal unary translations. Its
external output action remains exactly $f_O$, because the forest's output action
is $f_O$ as well. This theorem applies at arbitrary allowance and to arbitrary
operands admitting $f$, including operands with nested boxes. It does not itself
construct the separate box-body renaming correspondence needed for a lambda.

\lean{Copying/UnaryBoundaryRenaming}{32}{34}{UnaryBoundaryRenaming.unaryForestIso}
\lean{Copying/UnaryBoundaryRenaming}{36}{43}{UnaryBoundaryRenaming.unary_forest_compatible}
\lean{Copying/UnaryBoundaryRenaming}{46}{47}{UnaryBoundaryRenaming.unary_boundaryIso}
\lean{Copying/UnaryBoundaryRenaming}{50}{55}{UnaryBoundaryRenaming.unary_boundaryIso_output}

\subsection{Induction over the original syntax}
The unary-fragment predicate is an indexed predicate on the existing
\texttt{Term}: atomic terms belong to it, and applying any original unary
operation preserves membership. It is not a new term datatype or a replacement
translation. Transporting a graph along equality of allowances commutes with
the original unary clause; eliminate that equality to prove the statement.
Now induct on fragment membership. The leaf case is the previous atomic
renaming theorem. For a unary node, expose the original translation and renaming
clauses, move the allowance equality through the unary constructor, and apply
the constructor theorem to the induction hypothesis. The result has only the
fragment-membership and bijection hypotheses, with no machine, stability, or
template premise.

\lean{Copying/UnaryFragmentRenaming}{10}{12}{UnaryFragmentRenaming.UnaryFragment}
\lean{Copying/UnaryFragmentRenaming}{15}{18}{UnaryFragmentRenaming.castDepth_unary}
\lean{Copying/UnaryFragmentRenaming}{22}{32}{UnaryFragmentRenaming.unaryFragment_rename}

A list of unary operators can be folded around an original term. Induction
on the list shows that folding around an atomic term belongs to the fragment,
and the fragment theorem then proves renaming for arbitrarily long chains.
This supplies an unbounded syntactic family, rather than a fixed set of examples.

\lean{Copying/UnaryFragmentRenaming}{35}{37}{UnaryFragmentRenaming.unaryChain}
\lean{Copying/UnaryFragmentRenaming}{39}{43}{UnaryFragmentRenaming.unaryChain_fragment}
\lean{Copying/UnaryFragmentRenaming}{45}{50}{UnaryFragmentRenaming.unaryChain_rename}

The three public example constructions use the constructor theorem directly.
For $\mathtt{ref}\ 1$, use the original constant automorphism with both ambient
resource bijections. For dereferencing an atom, use the original atom-renaming
isomorphism, retaining its instance edge. For dereferencing a variable, use the
original selected-variable isomorphism. All three have arbitrary finite variable
and atom scopes and preserve the original terms, labels, and sharing graphs.
They establish representation correspondence, not behavioural equivalence
between allocation and a constant result.

\lean{Copying/UnaryFragmentRenaming}{57}{59}{UnaryFragmentRenaming.referenceOne_boundaryRename}
\lean{Copying/UnaryFragmentRenaming}{62}{66}{UnaryFragmentRenaming.dereferenceAtom_boundaryRename}
\lean{Copying/UnaryFragmentRenaming}{69}{73}{UnaryFragmentRenaming.dereferenceVariable_boundaryRename}

\paragraph{Updated remaining work.}
Unary operations are removed from the remaining bijective translation-renaming
cases. Application, binary operations, and lambda boxes remain, together with
non-bijective renaming and graph-level substitution. The concrete copying UAM,
its representation laws, the eleven template families, semantic weakening,
and Figure~22's observational-equivalence chain are still required. The full
build and all three axiom audits pass for the present milestone without
\texttt{sorryAx} or project-specific axioms.

As of 25 October 2026, the estimated completion is approximately
\textbf{45\% by mathematical effort}, with a broad plausible range of
\textbf{30--55\%}. This supersedes the earlier estimate after the serial-only
milestone. The increase reflects a completed constructor and an unbounded
fragment theorem; it does not count infrastructure declarations as completed
machine or template proofs.


The application and binary obligations recorded at this milestone are subsequently
discharged in Section~\ref{sec:copying-lambda-free-renaming}; Appendix~B records
the current remaining tasks.
